\documentclass[10pt,a4paper]{amsart}
\usepackage[margin=19mm]{geometry}
\usepackage{amsmath,amssymb,mathtools}
\usepackage[scr=rsfs,cal=euler]{mathalfa}
\usepackage{xfrac}
\usepackage[unicode,hidelinks,bookmarks=false]{hyperref}
\newcommand{\ie}{i.e.\ }
\newcommand{\cf}{cf.\ }
\newcommand{\resp}{respectively\ }
\newcommand{\Subsection}{\subsection}
\providecommand{\nobreakdash}{\nobreak\hbox{-}}
\setlength{\emergencystretch}{2em}
\multlinegap0pt
\usepackage{amssymb}
\usepackage{mathtools}
\newtheorem{proposition}{Proposition}
\newcommand{\Ocoeff}{\mathcal O_E}
\title{Space of norms on locally algebraic representations}
\author{Alexandre Pyvovarov}
\date{Source: arXiv:2607.17605v1 (CC BY 4.0)}
\begin{document}
\maketitle

\noindent\emph{Source and licence.} Space of norms on locally algebraic representations. Version arXiv:2607.17605v1 (CC BY 4.0), distributed by arXiv under the Creative Commons Attribution 4.0 International licence, http://creativecommons.org/licenses/by/4.0/, as recorded in the arXiv licence metadata for this exact version and shown on the abstract page https://arxiv.org/abs/2607.17605v1. Full licence notice: \texttt{licenses/arxiv-2607.17605-CC-BY-4.0.txt}.\par\smallskip

\noindent\textit{Editorial scope.} Counted unit: the article's proposition prop:one-wall-Bernstein-Iwahori-positivity (source0.tex lines 3956--4140). The article's complete original proof of it is delivered with it (source0.tex lines 4142--4224, one \texttt{proof} environment of 83 lines). No mathematics is written, repaired or re-derived: the statement and the proof are the article's own lines, in the article's own notation, and no retained line is substituted or re-flowed.  Retained uncounted as the source-local context the proof needs: lines 2126--2129.  Nothing was omitted from any retained span, so the package carries no omission record.  No retained line contains a citation, so the wrapper carries no bibliography and no bibliography entry.  No retained line contains a figure environment, an image-inclusion command or drawing code, so no image is delivered and the package declares no asset dependency.  Editorial formatting only, in addition: an AMS \texttt{amsart} preamble that restates the article's own declarations and macros used by the retained lines, namely the environments \texttt{proposition} on separate counters, together with the standard packages those lines need.  The retained environments print in this document as Proposition 1 in order of appearance, whereas the article prints the same items with the numbering of its own full document.  The complete native source of the article as distributed in the arXiv e-print for this exact version is bundled unchanged as \texttt{sources/205584/source0.tex}.
\begin{equation}\label{eq:matrix-valuation-definition}
 v_E(U)=\min_{a,b}v_E(u_{ab}),
 \qquad v_E(0)=+\infty.
\end{equation}

\begin{proposition}[Explicit one-wall Bernstein--Iwahori positivity]
\label{prop:one-wall-Bernstein-Iwahori-positivity}
Put
\[
 h_j:=\langle\eta_j,\alpha_j\rangle\in\mathbb Z_{\geq0},
\]
so that the rank-one Bernstein relation at the \(j\)-th wall is
\begin{equation}\label{eq:explicit-one-wall-Bernstein-relation}
 T_{s_j}\Theta_{\eta_j}
 =
 \Theta_{s_j\eta_j}T_{s_j}
 +(q_\rho-1)
 \sum_{r=0}^{h_j-1}
 \Theta_{\eta_j-r\alpha_j^\vee}.
\end{equation}
For the principal branch set
\[
 \zeta_{j,\mathrm{pr}}:=s_j\eta_j,
\]
and for the \(r\)-th correction branch, \(0\leq r<h_j\), set
\[
 \zeta_{j,r}:=\eta_j-r\alpha_j^\vee.
\]
Assume the positive-prefix decompositions
\begin{align}
 (\zeta_{j,\mathrm{pr}})_{\mathrm{nc}}
 &=
 \sum_{k=1}^{\ell-1}
 m_{j,\mathrm{pr},k}\omega_{kd}^\vee,
 &
 m_{j,\mathrm{pr},k}&\in\mathbb Z_{\geq0},
 \label{eq:one-wall-principal-prefix-explicit}\\
 (\zeta_{j,r})_{\mathrm{nc}}
 &=
 \sum_{k=1}^{\ell-1}
 m_{j,r,k}\omega_{kd}^\vee,
 &
 m_{j,r,k}&\in\mathbb Z_{\geq0},
 \qquad 0\leq r<h_j,
 \label{eq:one-wall-correction-prefix-explicit}
\end{align}
and suppose
\begin{equation}\label{eq:one-wall-defects-nonnegative}
 \Delta_k\geq0
 \qquad(1\leq k<\ell).
\end{equation}

Fix the \(\Ocoeff\)-lattices in all finite-dimensional type spaces that
underlie the chosen type--Bruhat normal form.  Let
\[
 U_{j,\mathrm{pr}}
 \quad\text{and}\quad
 U_{j,r}\ \ (0\leq r<h_j)
\]
be the type-intertwining matrices attached respectively to the principal
branch and to the correction branches of
\eqref{eq:explicit-one-wall-Bernstein-relation}.  The hypothesis that the
type intertwiners are \emph{integral} means explicitly that, for the source
and target type lattices \(\Lambda_{\mathrm{in}}\) and
\(\Lambda_{\mathrm{out}}\) of each branch,
\begin{equation}\label{eq:one-wall-type-intertwiner-lattice-integrality}
 U(\Lambda_{\mathrm{in}})
 \subseteq
 \Lambda_{\mathrm{out}}.
\end{equation}
Equivalently, in the fixed integral bases,
\begin{equation}\label{eq:one-wall-type-intertwiner-matrix-integrality}
 U_{j,\mathrm{pr}}\in M_{r_{\mathrm{pr}}}(\Ocoeff),
 \qquad
 U_{j,r}\in M_{r_r}(\Ocoeff),
\end{equation}
with the evident rectangular variant when the source and target type blocks
have different ranks.  In terms of the matrix valuation
\eqref{eq:matrix-valuation-definition}, this is the pair of inequalities
\begin{equation}\label{eq:one-wall-type-intertwiner-valuation-integrality}
 v_E(U_{j,\mathrm{pr}})\geq0,
 \qquad
 v_E(U_{j,r})\geq0
 \quad(0\leq r<h_j).
\end{equation}

After the segment sign character is applied, write the complete normalized
coefficient matrix of the principal branch in the form
\begin{equation}\label{eq:one-wall-principal-complete-coefficient}
 M_{j,\mathrm{pr}}
 =
 u_{j,\mathrm{pr}}\,
 q_\rho^{\,c_{j,\mathrm{pr}}}
 (-1)\,
 \chi_{\mathfrak m}
 \!\left(\Theta_{\zeta_{j,\mathrm{pr}}}\right)
 \eta_{\mathrm{alg}}
 \!\left(\zeta_{j,\mathrm{pr}}\right)
 U_{j,\mathrm{pr}},
\end{equation}
and the complete normalized coefficient matrix of the \(r\)-th correction
branch in the form
\begin{equation}\label{eq:one-wall-correction-complete-coefficient}
 M_{j,r}
 =
 u_{j,r}\,
 q_\rho^{\,c_{j,r}}
 (q_\rho-1)\,
 \chi_{\mathfrak m}
 \!\left(\Theta_{\zeta_{j,r}}\right)
 \eta_{\mathrm{alg}}
 \!\left(\zeta_{j,r}\right)
 U_{j,r}.
\end{equation}
Here
\[
 u_{j,\mathrm{pr}},u_{j,r}\in\Ocoeff^\times,
 \qquad
 c_{j,\mathrm{pr}},c_{j,r}\in\mathbb Z_{\geq0},
\]
allow for the unit factors and nonnegative powers of the Hecke parameter
coming from the chosen local reduction convention.  Let $v_E(M)=\min_{a,b}v_E(M_{ab})$, where $M$ is a matrix. Then the principal branch satisfies the explicit identity
\begin{align}
 v_E(M_{j,\mathrm{pr}})
 &=
 c_{j,\mathrm{pr}}v_E(q_\rho)
 +
 \kappa_{E,F}
 \sum_{k=1}^{\ell-1}
 m_{j,\mathrm{pr},k}\Delta_k
 +
 v_E(U_{j,\mathrm{pr}})
 \notag\\
 &\geq0,
 \label{eq:one-wall-principal-valuation-nonnegative}
\end{align}
and, for every \(0\leq r<h_j\), the \(r\)-th correction branch satisfies
\begin{align}
 v_E(M_{j,r})
 &=
 c_{j,r}v_E(q_\rho)
 +
 v_E(q_\rho-1)
 +
 \kappa_{E,F}
 \sum_{k=1}^{\ell-1}
 m_{j,r,k}\Delta_k
 +
 v_E(U_{j,r})
 \notag\\
 &=
 c_{j,r}v_E(q_\rho)
 +
 \kappa_{E,F}
 \sum_{k=1}^{\ell-1}
 m_{j,r,k}\Delta_k
 +
 v_E(U_{j,r})
 \notag\\
 &\geq0.
 \label{eq:one-wall-correction-valuation-nonnegative}
\end{align}
In particular, every individual principal term and every individual
correction term in the rank-one expansion has nonnegative normalized
smooth--algebraic--type valuation.

The same conclusion holds recursively on every correction branch satisfying
the positive-prefix hypothesis.  More precisely, if a terminal term of a
recursive straightening is represented by a sequence of branch matrices
\[
 M_{\gamma}
 =
 M_{j_s,\varepsilon_s}\cdots
 M_{j_1,\varepsilon_1},
 \qquad
 \varepsilon_t\in
 \{\mathrm{pr}\}\cup
 \{0,\ldots,h_{j_t}-1\},
\]
then
\begin{equation}\label{eq:one-wall-recursive-product-valuation}
 v_E(M_\gamma)
 \geq
 \sum_{t=1}^s
 v_E(M_{j_t,\varepsilon_t})
 \geq0.
\end{equation}
Consequently every terminal branch coefficient is integral, and any finite
sum of such terminal coefficients is integral as well.
\end{proposition}

\begin{proof}
For every noncentral translation
\[
 \zeta_{\mathrm{nc}}
 =
 \sum_{k=1}^{\ell-1}m_k\omega_{kd}^\vee
 \qquad(m_k\geq0),
\]
the valuation identity established above gives
\begin{equation}\label{eq:one-wall-translation-valuation}
 v_E\!\left(
 \chi_{\mathfrak m}(\Theta_\zeta)
 \eta_{\mathrm{alg}}(\zeta)
 \right)
 =
 \kappa_{E,F}
 \sum_{k=1}^{\ell-1}m_k\Delta_k.
\end{equation}
The central part contributes valuation \(0\), because the normalized
smooth--algebraic central character is unitary by weak admissibility.
Applying \eqref{eq:one-wall-translation-valuation} to
\(\zeta_{j,\mathrm{pr}}\) and \(\zeta_{j,r}\) gives the terms appearing
in \eqref{eq:one-wall-principal-valuation-nonnegative} and
\eqref{eq:one-wall-correction-valuation-nonnegative}.

Next,
\[
 v_E(u_{j,\mathrm{pr}})
 =
 v_E(u_{j,r})
 =
 v_E(-1)
 =
 v_E(q_\rho-1)
 =
 0,
\]
because the \(u\)'s and \(q_\rho-1\) are units in \(\Ocoeff\).  Moreover,
\[
 c_{j,\mathrm{pr}}v_E(q_\rho)\geq0,
 \qquad
 c_{j,r}v_E(q_\rho)\geq0,
\]
since \(c_{j,\mathrm{pr}},c_{j,r}\geq0\), while
\eqref{eq:one-wall-type-intertwiner-valuation-integrality} gives
\[
 v_E(U_{j,\mathrm{pr}})\geq0,
 \qquad
 v_E(U_{j,r})\geq0.
\]
Finally, by
\eqref{eq:one-wall-principal-prefix-explicit},
\eqref{eq:one-wall-correction-prefix-explicit}, and
\eqref{eq:one-wall-defects-nonnegative},
\[
 \sum_{k=1}^{\ell-1}
 m_{j,\mathrm{pr},k}\Delta_k\geq0,
 \qquad
 \sum_{k=1}^{\ell-1}
 m_{j,r,k}\Delta_k\geq0.
\]
Adding these contributions proves
\eqref{eq:one-wall-principal-valuation-nonnegative} and
\eqref{eq:one-wall-correction-valuation-nonnegative}.

For a recursively produced terminal branch, the matrix valuation is
submultiplicative:
\[
 v_E(AB)\geq v_E(A)+v_E(B).
\]
Iterating this inequality gives
\eqref{eq:one-wall-recursive-product-valuation}.  If several terminal
branches contribute to the same matrix coefficient, the ultrametric
inequality gives
\[
 v_E\!\left(\sum_\gamma M_\gamma\right)
 \geq
 \min_\gamma v_E(M_\gamma)
 \geq0.
\]
Thus neither recursive multiplication nor collection of branches can
destroy integrality.
\end{proof}

\end{document}
